% Round 17. NF217. New separate native refinement of CD217.
% Does not change the already reviewed CD217 theorem or its published input.
\section{Native preparation of degenerate quadratic points and a field-degree bound}
\label{r17nf:section}

\paragraph{Purpose and inputs.}
CD217 proved the even-dimensional residue-two case using the
published Jannsen--Saito henselian normal form, which permits a
finite etale extension of the valuation ring. Here we give a
native construction over the completed valuation ring of the
original node residue field. The only preparation theorem used
is the formal one-variable Weierstrass theorem over a complete
DVR. A faithful-flatness argument then proves that its unit is
an actual regular function. No formal coordinate change is
assumed to algebraize.

The published preparation input is, for example, the classical
specialization of Laurent Berger, \emph{The Weierstrass preparation
theorem and resultants of p-adic power series}, Section 1 and
Theorem B, \url{https://arxiv.org/pdf/1910.05319} (published in
\emph{Muenster Journal of Mathematics} 14 (2021), 155--163).
The complete-DVR statement has been checked in that primary
source. The algebraic isometry slice is the independently
proved C217 lemma \texttt{r17c2:isometry-slice}.
No novelty claim is made for a Morse or Weierstrass normal-form
principle.

\begin{lemma}[The prepared unit is regular on an etale curve]
\label{r17nf:regular-unit}
Let $R$ be a complete DVR, and let $(C,p)$ be a pointed
smooth $R$-curve equipped with a morphism
$z:C\to\mathbf A^1_R$ which is etale at $p$, sends $p$
to the closed origin, and induces the same residue field
there. Put $A=\mathcal O_{C,p}$. If $g\in A$ has
reduction of order exactly $s\ge1$ in the local parameter
$z$ on the special fiber, then there are a monic polynomial
\[
        T(z)=z^s+b_{s-1}z^{s-1}+\cdots+b_0\in R[z],
                       \qquad b_i\in\mathfrak m_R,
\]
and a unit $u\in A^\times$ such that $g=uT(z)$ in $A$.
\end{lemma}

\begin{proof}
The given pointed etale coordinates identify the completed
noetherian local ring with
\[
                             \widehat A\simeq R[[z]].
\]
Here completeness of $R$ is essential. The formal
Weierstrass theorem gives $\widehat u\in R[[z]]^\times$
and the displayed distinguished polynomial $T\in R[z]$
with $g=\widehat uT$ in $\widehat A$.
The polynomial already belongs to $A$. Completion is
faithfully flat, so contraction of its principal ideal gives
\[
                         (T)\widehat A\cap A=(T)A.
 \tag{NF217.1}\label{r17nf:ideal-contraction}
\]
Consequently $g=uT$ for an actual $u\in A$.
The ring $\widehat A=R[[z]]$ is a domain and $T\ne0$,
so this $u$ maps to $\widehat u$. Its residue is nonzero,
hence $u$ is a unit of the local ring $A$. After shrinking
around $p$ it is a regular invertible function on $C$.
No descent of an arbitrary formal series or formal
isomorphism is being asserted: only divisibility by the
finite polynomial $T$, already in $A$, is descended.
\end{proof}

\begin{proposition}[Native quadratic preparation in even relative dimension at two]
\label{r17nf:normal-form}
Let $R=K^\circ$ for a finite $K/\mathbf Q_2$, and let
$\mathcal X/R$ be flat of finite type with smooth generic
fiber. Let $p$ be an ordinary quadratic special-fiber point
of even relative dimension $n\ge2$. Write $k_p=\kappa(p)$
and let $R_p/R$ be the unramified extension with residue
field $k_p$.

There is a common residue-isomorphic pointed etale
neighborhood of $\mathcal X$ at $p$ and of
\[
 \operatorname{Spec}R_p[y_1,\ldots,y_n,z]
                   /(P(y)+z^2+bz+c),
 \tag{NF217.2}\label{r17nf:equation}
\]
where $P$ is a quadratic form over $R_p$ with unimodular
polar matrix, $b,c\in\mathfrak m_{R_p}$, and
\[
                          0\ne b^2-4c\in\mathfrak m_{R_p}.
 \tag{NF217.3}\label{r17nf:discriminant}
\]
No further extension of the node residue field is needed
for this normal form. All coordinates and units in the
construction are regular on actual etale neighborhoods.
\end{proposition}

\begin{proof}
Work over $R_p$ at the diagonal point, which is rational
on the special fiber; write $R$ for this ring in the
local calculation. This local unramified chart is not
a replacement of the global model field.
The smooth-ambient lemma of ON17, whose Tor and
Jacobian proof does not use $2^{-1}$, writes
$\mathcal X=V(f)$ in a smooth ambient $W/R$ of relative
dimension $n+1$ near the selected point.

\emph{Choice of the transverse variables.}
The special quadratic form has a polar radical of
dimension one. Indeed its alternating polar matrix
has odd size, and a larger radical would contain a
nonzero zero of the restricted square over the algebraic
closure, contradicting smoothness of the projectivized
tangent quadric. Its value on the radical is nonzero
by the same smoothness argument.
Choose a native complement on which the polar form is
nondegenerate, and lift its basis together with a radical
basis vector to relative etale coordinates
\[
                              (x_1,\ldots,x_n,z)
\]
on $W$, all zero at the special point. The special
quadratic term then has the form
$\overline P(x)+\alpha z^2$ with
$\alpha\ne0$ and invertible polar matrix on the $x$
variables. This is a decomposition over the native
residue field, not a splitting into hyperbolic planes.

\emph{The critical curve and the exact transverse square.}
Express
$df=\sum_i F_i\,dx_i+F_z\,dz$.
The scheme $C=V(F_1,\ldots,F_n)$ is smooth of relative
dimension one and $z:C\to\mathbf A^1_R$ is etale near
the selected point, because the $x$-Jacobian is the
invertible transverse polar block. Shrink to that
neighborhood. Henselian lifting of its closed point
in the fiber $z=0$ gives a native section
$a:\operatorname{Spec}R\to C$ with $z(a)=0$.

Form the actual etale base change
\[
                         W_1=W\times_{\mathbf A^1_R}C.
\]
It has the diagonal section $\sigma:C\to W_1$ over $C$.
Let $u_i=x_i-x_i|_C$ be its relative coordinates, and
put $g=f|_C$. On a pointed neighborhood the section
ideal is $(u_1,\ldots,u_n)$. The relative differential
of $f$ over $C$ vanishes on this section by definition
of $C$. The conormal sequence therefore gives the exact
ideal membership
\[
                    f-g\in(u_1,\ldots,u_n)^2.
\]
Choose regular coefficients and write
\[
               f-g=\sum_{i\le j}A_{ij}u_iu_j=q_w(u).
 \tag{NF217.4}\label{r17nf:transverse-form}
\]
Along the native $R$-section $\sigma a$ its coefficient
form is a quadratic form $q_0$ over $R$ with unimodular
polar matrix. The latter assertion follows from the
transverse Hessian modulo the maximal ideal.

Now apply the original C217 isometry-slice lemma with
base DVR $R$, smooth $R$-scheme $W_1$, and $R$-section
$\sigma a$. There is no need to invoke an unstated
version over an arbitrary base $C$. It gives a pointed
etale $W_2\to W_1$ and a regular matrix $M$, equal to
the identity along the lifted $R$-section, such that
$q_0(Mv)=q_w(v)$. Set $y=Mu$. Thus, exactly on this
actual neighborhood,
\[
                                f=q_0(y)+g.
 \tag{NF217.5}\label{r17nf:split-function}
\]
The morphism $(y,C):W_2\to\mathbf A^n_C$ is etale
at the selected point: relatively over $C$, one has
$dy=du$ there because $u=0$ and $M=1$.
The morphism $W_1\to C$ was smooth of relative dimension
$n$, so this is the ordinary relative Jacobian criterion.
We retain the actual function $g$ on $C$ and do not
replace its coordinate by a formal one.

\emph{One-variable preparation and absorption of its unit.}
Modulo the uniformizer, the equations $F_i=0$ have
linear part $\overline B_P x=0$. Thus in the completed
special critical curve each $x_i$ has order at least
two in $z$. Substituting into the special quadratic
term shows
\[
                            \overline g=\alpha z^2+O(z^3),
                                \qquad\alpha\ne0.
\]
Apply Lemma~\ref{r17nf:regular-unit} to the actual local
ring $\mathcal O_{C,p}$. It gives
\[
                  g=u(z)\bigl(z^2+bz+c\bigr),
       \qquad b,c\in\mathfrak m_R,\quad
                  u\in\mathcal O_{C,p}^{\times}.
 \tag{NF217.6}\label{r17nf:prepared}
\]
The notation $u(z)$ means a regular function on the
etale curve $C$, not necessarily a polynomial or a
function on the whole affine line. Shrink so that it
and its inverse are regular. Its value $u(a)$ is a
native unit of $R$.

Multiplying the hypersurface equation by this regular
unit's inverse yields
\[
                  u^{-1}f=u^{-1}q_0(y)+z^2+bz+c.
\]
The variable-coefficient quadratic form $u^{-1}q_0$
restricts along the retained native $R$-section to
$P=u(a)^{-1}q_0$, whose polar matrix is unimodular.
A second application of the same C217 slice, again
over $R$ and with that $R$-section, gives a pointed
etale neighborhood and a regular matrix $N$ with
$P(Nv)=u^{-1}q_0(v)$. Put $w=Ny$. Its relative
Jacobian over $C$ is invertible at the selected point.
Composing with the etale coordinate $z:C\to\mathbf A^1_R$
shows that $(w,z)$ is an actual etale map to
$\mathbf A^{n+1}_R$, and
\[
                          u^{-1}f=P(w)+z^2+bz+c
 \tag{NF217.7}\label{r17nf:exact-identity}
\]
as an identity of regular functions. Restricting this
map to $f=0$ gives a common pointed etale neighborhood
with \eqref{r17nf:equation}, without changing the
selected residue field.

\emph{The nonzero discriminant.}
Certainly $b^2-4c$ lies in the maximal ideal. Suppose
it were zero. If $b=c=0$, put $z_*=0$; otherwise
$b^2=4c$ with $c\in\mathfrak m_R$ gives
\[
                         v(b)=v(2)+\tfrac12v(c)>v(2).
\]
Hence $z_*=-b/2$ lies in $\mathfrak m_R$ in either
case. The target section $(w,z)=(0,z_*)$ has the same
closed point as the chosen pointed etale map. Since
$R$ is henselian, that map lifts the section to the
chosen neighborhood. The section lies on the
hypersurface and its generic point is singular:
all derivatives in $w$ vanish at zero and
$2z_*+b=0$. Etale morphisms preserve the smoothness
criterion, so this contradicts the assumed smooth
generic fiber of $\mathcal X$. Thus the discriminant
is nonzero, proving \eqref{r17nf:discriminant}.
\end{proof}

\begin{corollary}[An explicit simultaneous auxiliary-degree bound]
\label{r17nf:degree}
Let $\mathcal X/K^\circ$ be an integral flat projective
model as in the even-dimensional ordinary quadratic
case of CD217, with finitely many native closed nodes
$p_1,\ldots,p_t$. Put
\[
 d_i=[\kappa(p_i):k],\qquad
 f=\operatorname{lcm}(d_1,\ldots,d_t),\qquad
 s=\sum_{i=1}^t d_i.
\]
If there are no nodes take $f=1,s=0$ and $L=K$.
Otherwise there is a finite extension $L/K$ on which
the native semistable chains of CD217 exist and such
that
\[
                               [L:K]\le f\,2^s.
 \tag{NF217.8}\label{r17nf:bound}
\]
It can be taken Galois over $K$. Its full actual degree,
not the displayed upper bound, is the one used in every
original-field trace or obstruction descent.
\end{corollary}

\begin{proof}
For each native node, Proposition~\ref{r17nf:normal-form}
gives its quadratic polynomial over the unramified
extension $K_i/K$ of degree exactly $d_i$, with no
additional residue extension. Let $K_0/K$ be the
unramified extension of degree $f$. It contains each
$K_i$, and over its residue field each original point
has exactly $d_i$ rational conjugate points.

For each of the $d_i$ embeddings $K_i\to K_0$, include
the conjugate quadratic polynomial. There are at most
$s$ such polynomials in total. Each has nonzero
discriminant and has splitting field of degree at most
two over $K_0$. Their compositum $L$ consequently has
\[
                    [L:K_0]\le2^s,
                  \qquad [L:K]=f[L:K_0]\le f2^s.
\]
The chosen full set of conjugate polynomials is stable
under the cyclic group $\operatorname{Gal}(K_0/K)$.
Their compositum of splitting fields is therefore
normal, and it is separable in characteristic zero;
thus $L/K$ can be taken Galois without a further
unbounded normal closure.

Every quadratic root is integral with zero special
reduction. Translating by one root changes the linear
coefficient to $\beta$ with $\beta^2=b^2-4c\ne0$.
Its positive integral valuation over $L$ is the number
$q$ used in the CD217 chain construction. All special
points are now rational, so its native blowups give
one finite chain at each of them. The model remains
integral by flatness and geometric integrality of the
generic fiber. CD217 proves regularity, reduced strict
semistability and the all-characteristic quadric
boundary property of the resolved model.

The actual G conclusion and its exact specified
support are already those of CD217. This corollary
supplies a bounded choice of witness field; it does
not replace its normalization by a bound or by a
selected quadratic branch. The trace uses the exact
composite degree $[L:K]$ and all native residue factors.
\end{proof}

\paragraph{Limitations.}
The normal-form statement requires a complete native DVR;
its proof uses actual coefficients of a Weierstrass
polynomial in that ring. It does not assert the same
finite polynomial over an arbitrary noncomplete base.
The degree bound depends on the number and residue
degrees of the nodes, as it must for simultaneous
splitting of several quadratic characters. Neither
statement supplies ordinary quadratic reduction for
an arbitrary variety. The general strong conjecture
outside the proved reduction sectors remains open.
